\section{Experimental Setup}\label{sec:setup}

\textbf{Models and data.} We use two publicly available SST-2 classifiers from the Hugging Face hub,
BERT-Base with 12 encoder layers\footnote{\nolinkurl{textattack/bert-base-uncased-SST-2}} and DistilBERT
with 6 encoder layers\footnote{\nolinkurl{distilbert-base-uncased-finetuned-sst-2-english}}. All encoder
linear layers are factorised and mapped to crossbars; embeddings, layer normalisation, attention
score computation and the classification head (including BERT's pooler) stay digital. Accuracy is measured on all 872 sentences of the
SST-2 validation set, with inputs truncated to 128 tokens. Placement scores and Fisher information use
2{,}048 sentences from the SST-2 training set, the same sentences for every score.

\Cref{sec:mnli} repeats part of the study on MNLI with its own setup.

\textbf{Experiments.} \Cref{tab:experiments} lists the experiments. Hybrid placements keep 20\% of
the components in SLC unless the SLC share is the swept quantity. Each error source is varied on its own,
with all others disabled, so that its effect can be attributed. Where a configuration contains no random
error, such as IR drop or ADC resolution without noise, a single evaluation is exact and only one chip is
used.

\textbf{Baseline.} The unchanged HyFlexPIM hybrid: gradient placement, 20\% SLC, no mapping or
hardware-aware training. HyFlexPIM reports programming noise only; we also test it, without our methods,
under our other error sources and on MNLI.

\begin{table}[t]
\centering\footnotesize
\caption{Experiments. Chips: number of simulated chips per point.}\label{tab:experiments}
\setlength{\tabcolsep}{2.5pt}
\begin{tabular}{@{}lll c@{}}
\toprule
Experiment & Swept & Fixed & Chips \\
\midrule
Programming noise & $\sigma=0$--10\% & & 5 \\
SLC share & share 5--50\% & $\sigma=3,5,7.5\%$ & 20 \\
Stuck cells & $p=0.1$--5\% & & 20 \\
Read noise & $\sigma_r=0.5$--5\% & & 5 \\
Level noise & mid levels $\times1.5$, $\times2$ & $\sigma=5\%$ & 10 \\
Fault pattern & scattered, clusters, lines & $p=1\%$ & 10 \\
IR drop & $r=10^{-4}$--$3\times10^{-3}$ & & 1 \\
ADC resolution & 3--8 bits & & 1 \\
ADC + noise$^{*}$ & 4--8 bits & $\sigma=5\%$ & 3 \\
\midrule
Fault-aware mapping & $p=0.2$--2\% & & 10 \\
 & $p=0.2$--1\% & $\sigma=5\%$ & 10 \\
 & clusters, lines & $p=1\%$ & 10 \\
\midrule
Training & as mapping; $\sigma=0$--7.5\% & & 10 \\
\bottomrule
\end{tabular}

\smallskip
\raggedright$^{*}$DistilBERT only. ADC experiments use bit-serial inputs.
\end{table}

\textbf{Evaluation batches.} Crossbar inputs share one 8-bit scale per batch, so every placement in an
experiment uses the same batches: 128 sentences for the basic sweeps, mapping and training, and 64 for
the other non-idealities.

\textbf{Statistics.} We report the mean accuracy over chips and, where useful, the worst chip. Because
placements share chips (\cref{sec:model}), we compare two placements by their paired difference: the
number of correctly classified sentences, out of 872, by which one placement beats the other on the same
chip, or the same difference in accuracy points, averaged over chips and given with its standard error.

\textbf{Compute.} All runs use Kaggle notebooks with one NVIDIA T4 GPU, PyTorch and Hugging Face
Transformers, one self-contained script per stage. Code, self-tests and raw results are at
\url{https://github.com/Shivanshmittal023/reliable-reram-crossbar}.
